\chapter{Garis dan Segi Banyak}\label{ch:g4:geometry}

Dengan penggaris, penggaris siku dan jangka, geometri menjadi sebuah
keterampilan. Bab ini mengajarkan dua kedudukan istimewa garis ---
\omterm{def:g4:geometry:def}{tegak lurus} dan \omterm{def:g4:geometry:def}{sejajar} --- pada tataran alat, lalu berkeliling
mengenal keluarga \omterm{def:g4:geometry:polygon}{segi banyak}. (Teorinya menyusul pada
\cref{ch:g6:lines}.)

\section{Tegak lurus dan sejajar}

\begin{definition}[Tegak lurus, sejajar]\label{def:g4:geometry:def}
Dua garis disebut \emph{tegak lurus}\index{garis tegak lurus} jika
keduanya berpotongan membentuk \omterm{def:g3:shapes:rightangle}{sudut siku-siku} (diperiksa dengan
penggaris siku, ditandai dengan persegi kecil). Dua garis disebut
\emph{sejajar}\index{garis sejajar} jika keduanya menuju arah yang
sama dan tidak pernah berpotongan --- seperti dua rel kereta yang
jaraknya selalu tetap.
\end{definition}

\begin{method}[Menggambar dengan penggaris siku]\label{met:g4:geometry:draw}
\emph{\omterm{def:g4:geometry:def}{Garis tegak lurus} terhadap $d$ yang melalui titik $P$}:
letakkan satu tepi \omterm{def:g3:shapes:rightangle}{sudut siku-siku} itu sepanjang $d$, lalu geser
sampai tepi yang lain mencapai $P$, dan gambar.

\emph{\omterm{def:g4:geometry:def}{Garis sejajar} dengan $d$ yang melalui $P$}: letakkan tepi
penggaris siku sepanjang $d$, tekan penggaris biasa pada tepi yang
lain sebagai rel, geser penggaris siku sepanjang rel itu sampai ke
$P$, lalu gambar.
\end{method}

\begin{omfigure}
\begin{tikzpicture}[scale=0.85]
  \draw[thick] (-0.4,0) -- (4.2,0) node[below, font=\small] {$d$};
  \draw[omDef, thick] (1.8,-0.9) -- (1.8,2.3);
  \draw[thin] (1.8,0) ++(-0.25,0) -- ++(0,0.25) -- ++(0.25,0);
  \fill (1.8,1.6) circle (1.6pt) node[right, font=\small] {$P$};
  \node[below, font=\small] at (1.9,-1.1) {tegak lurus melalui $P$};
  \begin{scope}[xshift=7.4cm]
  \draw[thick] (-0.4,0) -- (4.2,0) node[below, font=\small] {$d$};
  \draw[omThm, thick] (-0.4,1.6) -- (4.2,1.6);
  \fill (1.6,1.6) circle (1.6pt) node[above, font=\small] {$P$};
  \draw[Stealth-Stealth, thin] (3.4,0) -- (3.4,1.6)
    node[midway, right, font=\small] {jaraknya sama di mana-mana};
  \node[below, font=\small] at (1.9,-1.1) {sejajar melalui $P$};
  \end{scope}
\end{tikzpicture}

{\small Kedua lukisan itu. \omterm{def:g4:geometry:def}{Garis sejajar} menjaga jarak yang tetap:
mengukur jaraknya di dua tempat adalah pemeriksaan yang baik.}
\end{omfigure}

\begin{example}[Mengenalinya dalam gambar]\label{ex:g4:geometry:spot}
Pada persegi panjang, setiap sisi \omterm{def:g4:geometry:def}{tegak lurus} terhadap dua sisi
tetangganya dan \omterm{def:g4:geometry:def}{sejajar} dengan sisi yang berhadapan. Kertas berpetak
adalah jaringan \omterm{def:g4:geometry:def}{garis sejajar} dan \omterm{def:g4:geometry:def}{tegak lurus} yang sudah jadi ---
itulah yang membuatnya begitu nyaman untuk menggambar.
\end{example}

\section{Segi banyak}

\begin{definition}[Segi banyak]\label{def:g4:geometry:polygon}
\emph{Segi banyak}\index{segi banyak} adalah bangun tertutup yang
tepinya tersusun dari ruas garis lurus, yaitu \emph{sisi}-sisinya;
titik pojoknya disebut \emph{titik sudut}\index{titik sudut}. Segi
banyak dinamai menurut banyaknya sisi:
\emph{segitiga} ($3$), \emph{segi empat}\index{segi empat} ($4$),
\emph{segi lima}\index{segi lima} ($5$),
\emph{segi enam}\index{segi enam} ($6$),
\emph{segi delapan}\index{segi delapan}
($8$).
\end{definition}

\begin{omfigure}
\begin{tikzpicture}[scale=0.62]
  \draw[thick, fill=omDef!10] (0,0) -- (2,0) -- (1.2,1.7) -- cycle;
  \node[below, font=\small] at (1,-0.3) {segitiga};
  \begin{scope}[xshift=3.4cm]
  \draw[thick, fill=omThm!10] (0,0) -- (2.2,0) -- (2.5,1.5) --
    (0.5,1.8) -- cycle;
  \node[below, font=\small] at (1.25,-0.3) {segi empat};
  \end{scope}
  \begin{scope}[xshift=7.4cm]
  \draw[thick, fill=omMeth!10] (90:1.1) -- (162:1.1) -- (234:1.1) --
    (306:1.1) -- (18:1.1) -- cycle;
  \node[below, font=\small] at (0,-1.4) {segi lima};
  \end{scope}
  \begin{scope}[xshift=10.9cm]
  \draw[thick, fill=omProp!10] (0:1.1) -- (60:1.1) -- (120:1.1) --
    (180:1.1) -- (240:1.1) -- (300:1.1) -- cycle;
  \node[below, font=\small] at (0,-1.4) {segi enam};
  \end{scope}
\end{tikzpicture}

{\small Keluarga \omterm{def:g4:geometry:polygon}{segi banyak} (\omterm{def:g4:geometry:polygon}{segi lima} dan \omterm{def:g4:geometry:polygon}{segi enam} yang digambar
di sini \emph{beraturan}: sisinya sama panjang dan sudutnya sama
besar). Lingkaran bukan \omterm{def:g4:geometry:polygon}{segi banyak}: tepinya tidak tersusun dari ruas garis.}
\end{omfigure}

\begin{example}[Menjelaskan sebuah segi banyak]\label{ex:g4:geometry:describe}
Permainan yang berguna: jelaskan sebuah gambar setepat mungkin
sehingga teman sekelasmu dapat menggambarnya ulang tanpa melihatnya.
``Sebuah \omterm{def:g4:geometry:polygon}{segi empat} $ABCD$ dengan $AB = 5$ cm, sudut di $A$
siku-siku, $AD = 3$ cm\,\dots'' --- kosakatalah yang membuat pesan
itu sampai.
\end{example}

\section{Lingkaran dan jangka}

\begin{definition}[Lingkaran]\label{def:g4:geometry:circle}
\emph{Lingkaran} berpusat $O$ dengan \omterm{def:g3:shapes:shapes}{jari-jari} $3$ cm adalah
himpunan semua titik yang berjarak tepat $3$ cm dari $O$. Jangka
menggambarnya: jarum di $O$, bukaan $3$ cm, lalu putar.
\emph{Diameter}\index{diameter} adalah ruas garis melalui pusat yang
menghubungkan dua titik pada lingkaran; panjangnya dua kali
\omterm{def:g3:shapes:shapes}{jari-jari}.
\end{definition}

\begin{example}[Jangka sebagai penjaga jarak]\label{ex:g4:geometry:compass}
Jangka bukan hanya menggambar lingkaran: ia \emph{memindahkan
jarak}. Untuk menandai semua titik yang berjarak $4$ cm dari titik
$A$, gambarlah lingkaran berpusat $A$ dengan \omterm{def:g3:shapes:shapes}{jari-jari} $4$ cm; untuk
membandingkan dua ruas garis tanpa penggaris, bukalah jangka pada
yang satu lalu ujilah pada yang lain.
\end{example}

\begin{method}[Menggambar ulang sebuah gambar di kertas berpetak]\label{met:g4:geometry:reproduce}
\begin{enumerate}
  \item Pilihlah satu \omterm{def:g4:geometry:polygon}{titik sudut} awal beserta letaknya di petak yang
        masih kosong;
  \item telusurilah gambar itu \omterm{def:g4:geometry:polygon}{titik sudut} demi \omterm{def:g4:geometry:polygon}{titik sudut} sambil
        membilang langkah petaknya (``$3$ ke kanan, $2$ ke atas'');
  \item hubungkan \omterm{def:g4:geometry:polygon}{titik sudutnya} berurutan, lalu bandingkan dengan
        modelnya: panjangnya sama, arahnya sama.
\end{enumerate}
\end{method}

\section{Latihan}

\begin{exercise}[$\star$]\label{exo:g4:geometry:1}
Gambarlah sebuah garis $d$ dan titik $P$ yang tidak terletak
padanya. Lukislah \omterm{def:g4:geometry:def}{garis tegak lurus} terhadap $d$ yang melalui $P$,
lalu \omterm{def:g4:geometry:def}{garis sejajar} dengan $d$ yang melalui $P$ (dengan cara rel
penggaris). Tandailah \omterm{def:g3:shapes:rightangle}{sudut siku-sikunya}.
\end{exercise}

\begin{exercise}[$\star$]\label{exo:g4:geometry:2}
Perhatikan huruf kapital E, H, N, T, Z. Pada masing-masing, carilah
sepasang goresan yang \omterm{def:g4:geometry:def}{sejajar} dan sepasang goresan yang \omterm{def:g4:geometry:def}{tegak lurus}
bila memang ada.
\end{exercise}

\begin{exercise}[$\star$]\label{exo:g4:geometry:3}
Sebutkan nama setiap \omterm{def:g4:geometry:polygon}{segi banyak}: $6$ sisi; $4$ sisi; $3$ sisi; $8$
sisi; $5$ sisi. Berapa \emph{\omterm{def:g4:geometry:polygon}{titik sudut}} yang dipunyai masing-masing?
\end{exercise}

\begin{exercise}[$\star$]\label{exo:g4:geometry:4}
Gambarlah sebuah \omterm{def:g4:geometry:polygon}{segi lima} (tidak harus beraturan) lalu berilah nama
\omterm{def:g4:geometry:polygon}{titik sudutnya} $A$, $B$, $C$, $D$, $E$. Berapa sisinya? Gambarlah
diagonalnya dari $A$ (ruas garis yang menghubungkan $A$ ke \omterm{def:g4:geometry:polygon}{titik
sudut} yang bukan tetangganya): ada berapa?
\end{exercise}

\begin{exercise}[$\star$]\label{exo:g4:geometry:5}
Gambarlah lingkaran berpusat $O$ dengan \omterm{def:g3:shapes:shapes}{jari-jari} $4$ cm; sebuah
\omterm{def:g4:geometry:circle}{diameter} $[AB]$; dan sebuah titik $M$ pada lingkaran yang tidak
terletak pada $[AB]$. Sebutkan panjang $OA$, $OM$ dan $AB$ tanpa mengukur.
\end{exercise}

\begin{exercise}[$\star$]\label{exo:g4:geometry:6}
Dengan jangka sebagai penjaga jarak
(\cref{ex:g4:geometry:compass}), bandingkan sisi sebuah segitiga
yang digambar teman sekelasmu lalu putuskan apakah segitiga itu punya dua sisi yang sama panjang.
\end{exercise}

\begin{exercise}[$\star$]\label{exo:g4:geometry:7}
Di kertas berpetak, gambarlah ulang jalur ini lalu tutuplah menjadi
\omterm{def:g4:geometry:polygon}{segi banyak}: mulailah dari titik $A$; $4$ ke kanan; $2$ ke atas; $3$
ke kiri; $2$ ke atas; lalu kembali ke $A$ dengan garis lurus. Apa
nama \omterm{def:g4:geometry:polygon}{segi banyak} yang kamu peroleh (bilanglah sisinya)?
\end{exercise}

\begin{exercise}[$\star$]\label{exo:g4:geometry:8}
Benar atau salah? Jelaskan dengan gambar.
``Dua garis yang \omterm{def:g4:geometry:def}{tegak lurus} selalu berpotongan.''
``Dua garis yang \omterm{def:g4:geometry:def}{sejajar} tidak pernah berpotongan.''
``Sebuah garis dapat \omterm{def:g4:geometry:def}{sejajar} dengan satu garis dan \omterm{def:g4:geometry:def}{tegak lurus}
terhadap garis yang lain.''
\end{exercise}

\begin{exercise}[$\star\star$]\label{exo:g4:geometry:9}
Gambarlah sebuah garis $d$, lalu garis $e$ yang \omterm{def:g4:geometry:def}{tegak lurus} terhadap
$d$, lalu garis $f$ yang \omterm{def:g4:geometry:def}{tegak lurus} terhadap $e$. Perhatikan $d$
dan $f$: apa yang kamu lihat? (Pengamatan ini menjadi teorema pada
\cref{ch:g6:lines}.)
\end{exercise}

\begin{exercise}[$\star\star$]\label{exo:g4:geometry:10}
Jelaskan gambar berikut dengan kata-kata kepada teman sekelasmu,
tanpa menunjukkannya: sebuah persegi dengan sisi $4$ cm, dengan
lingkaran berjari-jari $2$ cm yang berpusat di tengah persegi itu
(menyentuh keempat sisinya). Tulislah penjelasanmu sebagai langkah bernomor.
\end{exercise}

\begin{exercise}[$\star\star$]\label{exo:g4:geometry:11}
Berapa diagonal yang dipunyai \omterm{def:g4:geometry:polygon}{segi empat}? \omterm{def:g4:geometry:polygon}{Segi lima}? \omterm{def:g4:geometry:polygon}{Segi enam}?
Gambarlah ketiganya, bilanglah dengan teliti, lalu jelaskan pola
yang kamu lihat.

\end{exercise}
